\documentclass{article}
\usepackage[utf8]{inputenc}
\usepackage{amsmath}
\usepackage{amssymb}
\usepackage{amsthm}
\usepackage{graphicx}
\usepackage{parskip}

\newtheorem{proposition}{Proposition}
\newtheorem{claim}{Claim}
\newtheorem{remark}{Remark}

\title{An equally spaced construction for $k$-sum-free sets}
\author{Ferran Espu\~na}
\date{\today}

\begin{document}

\maketitle

We work in \(\mathbb T=\mathbb R/\mathbb Z\). We will build open sets
\[
    A\subset \mathbb T
\]
with
\[
    0\notin A+A-kA.
\]
Equivalently, there are no \(x,y,z\in A\) satisfying \(x+y=kz\) in
\(\mathbb T\).

Fix an integer \(k\geq 3\). Choose, for the moment, an integer \(m\geq 1\), a
number \(\delta>0\), and a point \(a\in\mathbb T\). Consider
\[
    A=\bigcup_{r=0}^{m-1} I_r,
    \qquad
    I_r=\left(a+\frac r k,\; a+\frac r k+\delta\right).
\]
The intervals are equally spaced by \(1/k\). We will now discover how large
\(\delta\) can be.

\begin{claim}\label{claim:disjoint}
If
\[
    \delta<\frac1k
    \qquad\text{and}\qquad
    \frac{m-1}{k}+\delta<1,
\]
then the intervals \(I_0,\ldots,I_{m-1}\) are disjoint in \(\mathbb T\).
\end{claim}

\begin{proof}
In the chosen order, consecutive starting points are separated by \(1/k\), so
\(\delta<1/k\) prevents two consecutive intervals from overlapping. The second
inequality says that the last interval does not reach the first interval after
going around the circle.
\end{proof}

We next impose the \(k\)-sum-free condition. For \(0\leq r,s,t\leq m-1\), the
open interval \(I_r+I_s-kI_t\) has left and right endpoints
\begin{align*}
    L_{rst}
    &=
    (2-k)a+\frac{r+s-kt}{k}-k\delta, \\
    R_{rst}
    &=
    (2-k)a+\frac{r+s-kt}{k}+2\delta.
\end{align*}
Since \(kt/k=t\) is an integer, the index \(t\) only shifts the interval by an
integer. We therefore remove this integer shift and define
\[
    L'_{r,s}=(2-k)a+\frac{r+s}{k}-k\delta,
    \qquad
    R'_{r,s}=(2-k)a+\frac{r+s}{k}+2\delta.
\]
Then
\[
    I_r+I_s-kI_t=(L'_{r,s}-t,\; R'_{r,s}-t).
\]
Thus it is enough to place all the reduced intervals
\((L'_{r,s},R'_{r,s})\) between the same two consecutive integers.

\begin{claim}\label{claim:reduced_intervals}
Assume
\[
    L'_{0,0}=(2-k)a-k\delta\in\mathbb Z
\]
and
\[
    R'_{m-1,m-1}-L'_{0,0}\leq 1.
\]
Then \(A\) is \(k\)-sum-free.
\end{claim}

\begin{proof}
The reduced endpoints are increasing with \(r+s\). Hence, for every
\(0\leq r,s\leq m-1\),
\[
    L'_{0,0}
    \leq
    L'_{r,s}
    <
    R'_{r,s}
    \leq
    R'_{m-1,m-1}
    \leq
    L'_{0,0}+1.
\]
Thus each reduced interval \((L'_{r,s},R'_{r,s})\) lies between the same two
consecutive integers. Shifting by the integer \(-t\), each interval
\[
    I_r+I_s-kI_t=(L'_{r,s}-t,\;R'_{r,s}-t)
\]
also lies between two consecutive integers. Since these intervals are open,
integer endpoints are not included. Therefore no interval \(I_r+I_s-kI_t\)
contains an integer, and \(0\notin A+A-kA\) in \(\mathbb T\).
\end{proof}

The condition \((2-k)a-k\delta\in\mathbb Z\) can always be achieved after
\(\delta\) is chosen. For example, we may take
\[
    a=\frac{k\delta}{2-k},
\]
which makes \(L'_{0,0}=0\). In particular, when \(\delta\) is rational this
choice gives rational endpoints. We are therefore left with the size constraint
coming from Claim~\ref{claim:reduced_intervals}. Directly computing,
\[
    R'_{m-1,m-1}-L'_{0,0}
    =
    \frac{2(m-1)}{k}+(k+2)\delta.
\]
To make the intervals as large as possible for this value of \(m\), we take
equality:
\[
    \frac{2(m-1)}{k}+(k+2)\delta=1.
\]
Solving for \(\delta\) gives
\[
    \delta
    =
    \frac{1-\frac{2(m-1)}{k}}{k+2}
    =
    \frac{k-2m+2}{k(k+2)}.
\]
This is positive precisely when
\[
    m<\frac{k+2}{2}.
\]

\begin{proposition}\label{prop:construction}
Let \(k\geq 3\), and let \(m\) be an integer with
\[
    1\leq m<\frac{k+2}{2}.
\]
Set
\[
    \delta=\frac{k-2m+2}{k(k+2)}.
\]
Then there is a choice of \(a\in\mathbb T\) for which
\[
    A=\bigcup_{r=0}^{m-1}
    \left(a+\frac r k,\; a+\frac r k+\delta\right)
\]
is \(k\)-sum-free. Its measure is
\[
    \mu(A)=m\delta
    =
    \frac{m(k-2m+2)}{k(k+2)}.
\]
\end{proposition}

\begin{proof}
Choose \(a=k\delta/(2-k)\), so that \(L'_{0,0}=0\). The defining choice of
\(\delta\) gives the second hypothesis of Claim~\ref{claim:reduced_intervals}.
Hence \(A\) is \(k\)-sum-free.

It remains only to check that the intervals are disjoint. Since \(m\geq 1\),
\[
    \delta
    =
    \frac{k-2m+2}{k(k+2)}
    \leq
    \frac{k}{k(k+2)}
    <
    \frac1k.
\]
Also,
\[
    \frac{m-1}{k}+\delta
    <
    \frac{m}{k}
    <
    \frac{k+2}{2k}
    <1,
\]
because \(k\geq 3\). By Claim~\ref{claim:disjoint}, the intervals are
disjoint. The measure is therefore the sum of their lengths, namely
\(m\delta\).
\end{proof}

Now we optimize over \(m\). Put
\[
    \varphi(x)=x(k-2x+2).
\]
The density from Proposition~\ref{prop:construction} is
\[
    \frac{\varphi(m)}{k(k+2)}.
\]
The function \(\varphi\) is a downward-opening parabola with roots \(0\) and
\((k+2)/2\), so its real maximum occurs at
\[
    x_0=\frac{k+2}{4}.
\]
Therefore we choose an integer \(m_0\) nearest to \(x_0\). Then
\[
    \left|m_0-\frac{k+2}{4}\right|\leq \frac12.
\]
For \(k\geq 3\), this integer is allowed: indeed \(m_0\geq 1\), and
\[
    m_0
    \leq
    \frac{k+2}{4}+\frac12
    =
    \frac{k+4}{4}
    <
    \frac{k+2}{2}.
\]

\begin{claim}\label{claim:nearest_integer}
For this choice of \(m_0\),
\[
    \varphi(m_0)
    =
    \left\lfloor \frac{(k+2)^2}{8}\right\rfloor.
\]
\end{claim}

\begin{proof}
Completing the square gives
\[
    \varphi(x)
    =
    \frac{(k+2)^2}{8}
    -
    2\left(x-\frac{k+2}{4}\right)^2.
\]
Thus
\[
    \varphi\left(\frac{k+2}{4}\right)-\varphi(m_0)
    =
    2\left(m_0-\frac{k+2}{4}\right)^2
    \leq
    \frac12
    <1.
\]
Since \(\varphi(m_0)\) is an integer and
\(\varphi(m_0)\leq (k+2)^2/8\), this forces
\[
    \varphi(m_0)
    =
    \left\lfloor \frac{(k+2)^2}{8}\right\rfloor.
\]
\end{proof}

\begin{proposition}\label{prop:optimized_density}
For every \(k\geq 3\), the construction above gives a \(k\)-sum-free open set
of measure
\[
    \frac{\left\lfloor (k+2)^2/8\right\rfloor}{k(k+2)}.
\]
In particular this measure is always strictly larger than \(1/8\).
\end{proposition}

\begin{proof}
The formula follows from Proposition~\ref{prop:construction} and
Claim~\ref{claim:nearest_integer}. To prove that this is larger than \(1/8\),
use the estimate from the proof of Claim~\ref{claim:nearest_integer}:
\[
    \varphi(m_0)>
    \frac{(k+2)^2}{8}-1.
\]
Therefore
\[
    \mu(A)
    >
    \frac{\frac{(k+2)^2}{8}-1}{k(k+2)}.
\]
The right-hand side is larger than \(1/8\) exactly when
\[
    \frac{(k+2)^2}{8}-1>\frac{k(k+2)}{8},
\]
or equivalently
\[
    (k+2)^2-8>k(k+2).
\]
This reduces to \(2k-4>0\), which holds for every \(k\geq 3\).
\end{proof}

\begin{table}[ht]
\centering
\small
\caption{Optimized density of the equally spaced construction.}
\resizebox{\textwidth}{!}{%
\begin{tabular}{c|ccccccccccccc}
\(k\) & 3 & 4 & 5 & 6 & 7 & 8 & 9 & 10 & 11 & 12 & 13 & 14 & 15 \\
\hline
fraction
& \(1/5\) & \(1/6\) & \(6/35\) & \(1/6\) & \(10/63\) & \(3/20\)
& \(5/33\) & \(3/20\) & \(21/143\) & \(1/7\) & \(28/195\)
& \(1/7\) & \(12/85\) \\
decimal
& \(0.200000\) & \(0.166667\) & \(0.171429\) & \(0.166667\)
& \(0.158730\) & \(0.150000\) & \(0.151515\) & \(0.150000\)
& \(0.146853\) & \(0.142857\) & \(0.143590\) & \(0.142857\)
& \(0.141176\)
\end{tabular}%
}
\end{table}

\begin{remark}
This construction is not always optimal among finite unions of intervals. For
example, when \(k=4\), the equally spaced construction gives density \(1/6\),
but one can add a tiny extra interval and obtain density
\[
    \frac{11}{64}
    =
    \frac16+\frac1{192}.
\]
Thus the construction above should be viewed as a robust general lower bound,
not as a sharp formula for every \(k\).
\end{remark}

\end{document}
